% GENERATED FILE. Edit canonical lessons, then run npm run generate:books.
% canonical-source-hash: fnv1a64:49ec99f7a347528c
% canonical-lessons: PA-C337-mangna, PA-C337-benti, PA-C337-thorha, PA-C337-zara, PA-C337-mangvauna

\chapter{Asking Nicely}
\label{ch:pa-a2-asking-nicely}
\hlchaptermodality{337}

\begin{chapteropening}
\textbf{By the end of this chapter:} \emph{I can ask politely for something I need.}

Complete the last lesson of chapter 337: I can ask politely for something I need.
\end{chapteropening}

\section[\texorpdfstring{\pa{ਮੰਗਣਾ} (maṅgṇā)}{ਮੰਗਣਾ (maṅgṇā)}]{\pa{ਮੰਗਣਾ} (maṅgṇā) — to ask for}

\label{lesson:PA-C337-mangna}



\begin{quote}
Before the new one: say the Punjabi for to take out, then the Punjabi for perhaps.
\end{quote}

\subsection*{You'll want to know: \pa{ਮੰਗਣਾ}}
\textbf{\pa{ਮੰਗਣਾ}} — \emph{maṅgṇā} — \textquotedblleft{}to ask for\textquotedblright{}.

\begin{quote}
\pa{ਮੈਂ} \pa{ਕਦੇ} \pa{ਪੈਸੇ} \pa{ਨਹੀਂ} \pa{ਮੰਗਦਾ} \pa{ਹਾਂ।}
\end{quote}

\begin{quote}
\emph{maiṁ kade paise nahīṁ maṅgdā hāṁ.}
\end{quote}

\begin{quote}
I never ask for money.
\end{quote}

Past (I …): \pa{ਮੈਂ} \pa{ਪਾਣੀ} \pa{ਮੰਗਿਆ}, \emph{maiṁ pāṇī maṅgiā}.

\begin{cousinweb}
\textbf{\pa{ਮੰਗਣਾ}} (\emph{maṅgṇā}) is inherited from Sanskrit \emph{mārgati}, \textquotedblleft{}seeks, requests\textquotedblright{}; CDIAL lists Punjabi \emph{maṅgṇā}, \textquotedblleft{}to ask for\textquotedblright{}, under it (CDIAL 10074).
\end{cousinweb}

\begin{grammarlens}[title={What you've built}]
One more word for asking and answering.
\end{grammarlens}

\subsection*{Your turn}
\begin{itemize}
\raggedright
  \item \emph{Say it:} \emph{maṅgṇā}
  \item \emph{Say it:} \emph{maṅgṇā}, once more
  \item \emph{Say it:} \emph{maiṁ kade paise nahīṁ maṅgdā hāṁ.} — I never ask for money.
  \item \emph{Recall:} say the Punjabi for to take out, then the Punjabi for perhaps, then say \emph{maṅgṇā} again
\end{itemize}

\begin{tcolorbox}[breakable,colback=teal!4,colframe=teal!35!black,title={Before you move on}]
What does \pa{ਮੰਗਣਾ} mean? (\textquotedblleft{}To ask for\textquotedblright{}.) Say it once more.
\end{tcolorbox}
\section[\texorpdfstring{\pa{ਬੇਨਤੀ} (bentī)}{ਬੇਨਤੀ (bentī)}]{\pa{ਬੇਨਤੀ} (bentī) — a request}

\label{lesson:PA-C337-benti}



\begin{quote}
Before the new one: say the Punjabi for perhaps, then the Punjabi for to ask for.
\end{quote}

\subsection*{You'll want to know: \pa{ਬੇਨਤੀ}}
\textbf{\pa{ਬੇਨਤੀ}} — \emph{bentī} — \textquotedblleft{}a request\textquotedblright{}.

\begin{quote}
\pa{ਕੀ} \pa{ਤੁਸੀਂ} \pa{ਮੇਰੀ} \pa{ਇੱਕ} \pa{ਬੇਨਤੀ} \pa{ਮੰਨੋਗੇ}?
\end{quote}

\begin{quote}
\emph{kī tusī̃ merī ikk bentī mannoge?}
\end{quote}

\begin{quote}
Will you grant me one request?
\end{quote}

A feminine noun.

\begin{cousinweb}
\textbf{\pa{ਬੇਨਤੀ}} (\emph{bentī}) is inherited from Sanskrit \emph{vijñapti}, \textquotedblleft{}information, a request to a superior\textquotedblright{}; CDIAL lists Punjabi \emph{beṇtī}, \textquotedblleft{}request\textquotedblright{}, under it (CDIAL 11706).
\end{cousinweb}

\begin{grammarlens}[title={What you've built}]
One more word for asking and answering.
\end{grammarlens}

\subsection*{Your turn}
\begin{itemize}
\raggedright
  \item \emph{Say it:} \emph{bentī}
  \item \emph{Say it:} \emph{bentī}, once more
  \item \emph{Say it:} \emph{kī tusī̃ merī ikk bentī mannoge?} — Will you grant me one request?
  \item \emph{Recall:} say the Punjabi for perhaps, then the Punjabi for to ask for, then say \emph{bentī} again
  \item \emph{Write it:} copy \textbf{\pa{ਬੇਨਤੀ}} once, checking the whole word
\end{itemize}

\begin{tcolorbox}[breakable,colback=teal!4,colframe=teal!35!black,title={Before you move on}]
What does \pa{ਬੇਨਤੀ} mean? (\textquotedblleft{}A request\textquotedblright{}.) Say it once more.
\end{tcolorbox}
\section[\texorpdfstring{\pa{ਥੋੜ੍ਹਾ} (thoṛhā)}{ਥੋੜ੍ਹਾ (thoṛhā)}]{\pa{ਥੋੜ੍ਹਾ} (thoṛhā) — a little; not much}

\label{lesson:PA-C337-thorha}



\begin{quote}
Before the new one: say the Punjabi for to ask for, then the Punjabi for a request.
\end{quote}

\subsection*{You'll want to know: \pa{ਥੋੜ੍ਹਾ}}
\textbf{\pa{ਥੋੜ੍ਹਾ}} — \emph{thoṛhā} — \textquotedblleft{}a little; not much\textquotedblright{}.

\begin{quote}
\pa{ਕੀ} \pa{ਮੈਨੂੰ} \pa{ਥੋੜ੍ਹਾ} \pa{ਪਾਣੀ} \pa{ਮਿਲੇਗਾ}?
\end{quote}

\begin{quote}
\emph{kī mainū̃ thoṛhā pāṇī milegā?}
\end{quote}

\begin{quote}
May I have a little water?
\end{quote}

\begin{cousinweb}
\textbf{\pa{ਥੋੜ੍ਹਾ}} (\emph{thoṛhā}) is inherited from Sanskrit \emph{stoka}, \textquotedblleft{}a drop\textquotedblright{} (as an adverb, \textquotedblleft{}a little\textquotedblright{}); CDIAL lists Punjabi \emph{thoṛā}, \textquotedblleft{}a little, few\textquotedblright{}, under it (CDIAL 13720).
\end{cousinweb}

\begin{grammarlens}[title={What you've built}]
One more word for asking and answering.
\end{grammarlens}

\subsection*{Your turn}
\begin{itemize}
\raggedright
  \item \emph{Say it:} \emph{thoṛhā}
  \item \emph{Say it:} \emph{thoṛhā}, once more
  \item \emph{Say it:} \emph{kī mainū̃ thoṛhā pāṇī milegā?} — May I have a little water?
  \item \emph{Recall:} say the Punjabi for to ask for, then the Punjabi for a request, then say \emph{thoṛhā} again
  \item \emph{Write it:} copy \textbf{\pa{ਥੋੜ੍ਹਾ}} once, checking the whole word
\end{itemize}

\begin{tcolorbox}[breakable,colback=teal!4,colframe=teal!35!black,title={Before you move on}]
What does \pa{ਥੋੜ੍ਹਾ} mean? (\textquotedblleft{}A little; not much\textquotedblright{}.) Say it once more.
\end{tcolorbox}
\section[\texorpdfstring{\pa{ਜ਼ਰਾ} (zarā)}{ਜ਼ਰਾ (zarā)}]{\pa{ਜ਼ਰਾ} (zarā) — a bit (to soften a request); a moment}

\label{lesson:PA-C337-zara}



\begin{quote}
Before the new one: say the Punjabi for a request, then the Punjabi for a little.
\end{quote}

\subsection*{You'll want to know: \pa{ਜ਼ਰਾ}}
\textbf{\pa{ਜ਼ਰਾ}} — \emph{zarā} — \textquotedblleft{}a bit (to soften a request); a moment\textquotedblright{}.

\begin{quote}
\pa{ਜ਼ਰਾ} \pa{ਰੁਕੋ} \pa{ਅਤੇ} \pa{ਹੁਣ} \pa{ਨਾ} \pa{ਬੋਲੋ।}
\end{quote}

\begin{quote}
\emph{zarā ruko ate huṇ nā bolo.}
\end{quote}

\begin{quote}
Wait a moment and don't speak now.
\end{quote}

\begin{cousinweb}
\textbf{\pa{ਜ਼ਰਾ}} (\emph{zarā}) is a loan, usually traced through Persian to Arabic \emph{dharra}, \textquotedblleft{}a tiny particle\textquotedblright{}.
\end{cousinweb}

\begin{grammarlens}[title={What you've built}]
One more word for asking and answering.
\end{grammarlens}

\subsection*{Your turn}
\begin{itemize}
\raggedright
  \item \emph{Say it:} \emph{zarā}
  \item \emph{Say it:} \emph{zarā}, once more
  \item \emph{Say it:} \emph{zarā ruko ate huṇ nā bolo.} — Wait a moment and don't speak now.
  \item \emph{Recall:} say the Punjabi for a request, then the Punjabi for a little, then say \emph{zarā} again
  \item \emph{Write it:} copy \textbf{\pa{ਜ਼ਰਾ}} once, checking the whole word
\end{itemize}

\begin{tcolorbox}[breakable,colback=teal!4,colframe=teal!35!black,title={Before you move on}]
What does \pa{ਜ਼ਰਾ} mean? (\textquotedblleft{}A bit (to soften a request); a moment\textquotedblright{}.) Say it once more.
\end{tcolorbox}
\section[\texorpdfstring{\pa{ਮੰਗਵਾਉਣਾ} (maṅgvāuṇā)}{ਮੰਗਵਾਉਣਾ (maṅgvāuṇā)}]{\pa{ਮੰਗਵਾਉਣਾ} (maṅgvāuṇā) — to send for; to order (goods)}

\label{lesson:PA-C337-mangvauna}



\begin{quote}
Before the new one: say the Punjabi for a little, then the Punjabi for a bit (to soften a request).
\end{quote}

\subsection*{You'll want to know: \pa{ਮੰਗਵਾਉਣਾ}}
\textbf{\pa{ਮੰਗਵਾਉਣਾ}} — \emph{maṅgvāuṇā} — \textquotedblleft{}to send for; to order (goods)\textquotedblright{}.

\begin{quote}
\pa{ਅਸੀਂ} \pa{ਬਾਜ਼ਾਰ} \pa{ਤੋਂ} \pa{ਸਾਮਾਨ} \pa{ਨਹੀਂ} \pa{ਮੰਗਵਾਉਂਦੇ} \pa{ਹਾਂ।}
\end{quote}

\begin{quote}
\emph{asīṁ bāzār ton sāmān nahīṁ maṅgvāunde hāṁ.}
\end{quote}

\begin{quote}
We don't order goods from the market.
\end{quote}

Past (I …): \pa{ਮੈਂ} \pa{ਬਾਜ਼ਾਰ} \pa{ਤੋਂ} \pa{ਚਾਹ} \pa{ਮੰਗਵਾਈ}, \emph{maiṁ bāzār ton chāh maṅgvāī}.

\begin{cousinweb}
\textbf{\pa{ਮੰਗਵਾਉਣਾ}} (\emph{maṅgvāuṇā}) is the causative partner of \emph{maṅgṇā}, \textquotedblleft{}to ask for\textquotedblright{}, which goes back to Sanskrit \emph{mārgati}, \textquotedblleft{}seeks, requests\textquotedblright{}; CDIAL lists Punjabi \emph{maṅgṇā} under it (CDIAL 10074).
\end{cousinweb}

\begin{grammarlens}[title={What you've built}]
One more word for asking and answering.
\end{grammarlens}

\subsection*{Your turn}
\begin{itemize}
\raggedright
  \item \emph{Say it:} \emph{maṅgvāuṇā}
  \item \emph{Say it:} \emph{maṅgvāuṇā}, once more
  \item \emph{Say it:} \emph{asīṁ bāzār ton sāmān nahīṁ maṅgvāunde hāṁ.} — We don't order goods from the market.
  \item \emph{Recall:} say the Punjabi for a little, then the Punjabi for a bit (to soften a request), then say \emph{maṅgvāuṇā} again
  \item \emph{Write it:} copy \textbf{\pa{ਮੰਗਵਾਉਣਾ}} once, checking the whole word
\end{itemize}

\begin{tcolorbox}[breakable,colback=teal!4,colframe=teal!35!black,title={Before you move on}]
What does \pa{ਮੰਗਵਾਉਣਾ} mean? (\textquotedblleft{}To send for; to order (goods)\textquotedblright{}.) Say it once more.
\end{tcolorbox}
